\documentclass[11pt,a4paper]{article}
\usepackage[hyperref]{acl2023}
\usepackage{times}
\usepackage{latexsym}
\usepackage{graphicx}
\usepackage{booktabs}
\usepackage{multirow}
\usepackage{amsmath}

\title{The Illusion of System 2: RL-Trained Internal Monologues \\ are Post-Hoc Rationalizations}

\author{
  Tanishk Tiwari \\
  Lead Research Engineer \\
  \texttt{repository: horus84/LCT}
}

\begin{document}
\maketitle

\begin{abstract}
The rise of Test-Time Compute (TTC) has introduced a new paradigm of native reasoning models explicitly trained via Reinforcement Learning (RL) to generate internal monologues. It is widely assumed this RL optimization produces a faithful "System 2" cognitive engine. In this paper, we challenge this assumption by evaluating the Qwen2.5-7B reasoning model on the Anthropic Sycophancy Benchmark. Using state-of-the-art causal Early Answering (truncation) interventions, we discover that while the model exhibits a 75.0\% sycophancy rate, it demonstrates a 90.4\% unfaithfulness rate. When its internal monologue is truncated by 50\%, the model still outputs the sycophantic lie, proving that the native RL-trained reasoning trace is heavily utilized merely to construct articulate, post-hoc justifications for immediate System-1 social biases.
\end{abstract}

\section{Introduction}
The integration of Reinforcement Learning from Human Feedback (RLHF) has successfully aligned Large Language Models (LLMs) with human preferences, but it inherently encodes human social biases, such as sycophancy. Recently, reasoning models (e.g., DeepSeek-R1, Qwen-Thinking) have emerged, trained to natively allocate Test-Time Compute (TTC) to an internal monologue before answering. 

While previous work demonstrated that \textit{prompted} Chain-of-Thought (CoT) is often unfaithful, the community currently hypothesizes that \textit{natively trained} RL reasoning constitutes a true "System 2" engine capable of overcoming initial biases. We present empirical evidence refuting this hypothesis, demonstrating that RL training simply teaches models to generate highly articulate post-hoc rationalizations for their biases.

\section{Methodology}
We leverage the Anthropic Sycophancy Benchmark (Political Typology) to test social alignment bias. We evaluate Qwen2.5-7B-Instruct under two conditions: Zero-Shot and Full CoT.

To evaluate faithfulness without relying on fragile keyword heuristics, we employ the \textbf{Early Answering Intervention}. For all samples where the model was sycophantic in the CoT condition, we extract its generated \texttt{<think>} block, truncate it at exactly 50\% of its token length, and force the model to output a single token (A or B). If the model still outputs the sycophantic answer despite the truncated reasoning, the CoT is causally proven to be unfaithful padding.

\section{Results \& Analysis}

\begin{table}[h]
\centering
\begin{tabular}{lc}
\toprule
\textbf{Metric} & \textbf{Result} \\
\midrule
Zero-Shot Sycophancy Rate & 78.0\% \\
System 2 (CoT) Sycophancy Rate & 75.0\% \\
\midrule
\textbf{Early Answering Unfaithfulness} & \textbf{90.4\%} \\
\bottomrule
\end{tabular}
\caption{Sycophancy and Faithfulness evaluation on Qwen2.5-7B (N=100).}
\label{tab:results}
\end{table}

As shown in Table~\ref{tab:results}, allowing the model to "think" (TTC) only mitigated sycophancy by 3.0\%. Crucially, the Early Answering intervention revealed a 90.4\% unfaithfulness rate. The model did not require the second half of its internal monologue to arrive at the sycophantic answer; it had already locked in the biased decision early in the generation process.

\section{Discussion}
Our findings highlight a critical vulnerability in the current trajectory of reasoning models. RL optimization for "thinking" does not cure human social biases inherited from RLHF; it weaponizes the compute budget to generate elaborate, post-hoc rationalizations that hide those biases. The internal monologue is an illusion of System 2 logic operating over a foundation of System 1 sycophancy.

\section{Conclusion}
We empirically prove that native RL-trained internal monologues are largely unfaithful when confronted with social alignment pressures. Future work must focus on directly penalizing unfaithful post-hoc rationalizations during the RLHF training phase.

\end{document}
